% ch-a3-01 —— 第 3 篇 Fargues–Scholze，PDF 138–149，印刷页 125–136
\chapter{Langlands 纲领与曲线上的丛模空间}
\ysrc{125}
\begin{center}
{\small\itshape Laurent Fargues \quad Peter Scholze}
\end{center}
\medskip

\noindent\textbf{摘要.} 本文综述作者关于局部 Langlands 对应几何化的工作。我们解释曲线上 $G$-丛的模叠 $\Bun_G$ 的几何、范畴 $\mathcal D_{\mathrm{lis}}(\Bun_G,\Lambda)$ 的结构，以及借助前述内容构造局部 Langlands 参数。最后我们解释范畴几何化猜想。
\footnote{2020 年数学主题分类：主分类 11R39；次分类 14G22、14D24。第一作者受 ANR 项目 ANR-19-CE40-0015「COLOSS」与 ERC 高级资助 742608 GeoLocLang 资助。第二作者受 DFG Leibniz 奖资助，并受 DFG「卓越战略」EXC-2047/1 390685813 资助。\textcopyright 2025 American Mathematical Society。}

\medskip
\noindent\textbf{目录.}
\begin{enumerate}[leftmargin=2.2em,itemsep=0.1em]
\item 引言
\item Artin v-叠 $\Bun_G$
\item 范畴 $\mathcal D_{\mathrm{lis}}(\Bun_G,\Lambda)$
\item 几何 Satake 对应（[9, 第 VII 章]）
\item Langlands 参数与谱作用
\item 范畴几何化猜想
\item 一些最后的思考
\item 参考文献
\end{enumerate}

\section{引言}

固定局部域 $E$：其剩余域 $\F_q$ 有限、特征为 $p$，并取单值化元 $\pi$。于是或者 $E\cong\F_q((\pi))$（等特征情形），或者 $[E:\Q_p]<+\infty$（不等特征情形）。固定 $\F_q$ 的一个代数闭包 $\overline{\F}_q$，并把 $E$ 的相应最大非分歧扩张的完备化记作 $\breve E$；其 Frobenius 记作 $\sigma$。再固定一个素数 $\ell\neq p$。

设 $G$ 是 $E$ 上的约化群。对 $\Lambda\in\{\overline{\F}_\ell,\overline{\Q}_\ell\}$，我们说明如何构造半单局部 Langlands 对应
\[
\pi\longmapsto\varphi_\pi,
\]
它从 $G(E)$ 的取值于 $\Lambda$ 的不可约光滑表示出发，指向半单 Langlands 参数。

\ysrc{126}我们进一步说明如何把它「族化」到 $\Z_\ell$ 上，即构造为两个范畴中心之间的态射。为此我们使用几何 Langlands 纲领的方法，背景是 Fargues–Fontaine 曲线上 $G$-丛的模叠（[9]）。

\section{Artin v-叠 $\Bun_G$}

\subsection{定义与基本性质}

\subsubsection{相对曲线（[9, 第 II.1 章]）}

把 $\F_q$-完满 空间的范畴记作 $\Perf_{\F_q}$，并赋予它 \emph{v-拓扑}（[23]）。按定义，一族 完满 空间的态射 $(S_i\to S)_{i\in I}$ 是 v-覆盖，是指：对 $S$ 的任一拟紧开子集 $U$，存在 $I$ 的有限子集 $J$，以及 $S_j$（$j\in J$）的拟紧开子集 $V_j$，使得 $U=\bigcup_{j\in J}\mathrm{Im}(V_j\to S)$。这是概形的 fpqc 拓扑在 完满 空间上的类比。我们还将经常使用 完满 空间上的另一个 Grothendieck 拓扑：\emph{pro-étale 拓扑}；v-拓扑比 pro-étale 拓扑更细。为此回忆：与诺特仿射 adic 空间的范畴不同，仿射 完满 空间的范畴是完备的：对仿射 完满 空间的反向系统 $(\Spa(R_i,R_i^+))_i$，
\[
\varprojlim_i\Spa(R_i,R_i^+)=\Spa(R_\infty,R_\infty^+),
\]
其中若 $\varpi$ 是某个 $R_{i_0}$ 的伪单值化元在 $\varinjlim_i R_i^+$ 中的像，则 $R_\infty^+$ 是 $\varinjlim_i R_i^+$ 的 $\varpi$-进完备化，而 $R_\infty=R_\infty^+[\frac{1}{\varpi}]$。按定义，完满 空间的态射是 pro-étale 的，是指它在源与靶上都局部地可以写成以 étale 过渡态射连接的仿射 完满 空间的极限。pro-étale 拓扑是这样一种拓扑：其覆盖的定义方式与 v-拓扑相同，但要求族 $(S_i\to S)_{i\in I}$ 中的态射都是 pro-étale 的（必须注意，与 étale 态射不同，pro-étale 态射一般不是开映射；并且与 v-拓扑的情形一样，仅要求 $\coprod_{i\in I}|S_i|\to|S|$ 满射，不足以定义 pro-étale 覆盖）。

下文把
\[
*=\Spa(\overline{\F}_q)
\]
记作 $\Perf_{\F_q}$ 上 v-层拓扑斯的终对象。它不能由 完满 空间表示，但将是我们工作的基点。有时我们把情形底变到 $\Spa(C)$，其中 $C|\F_q$ 是代数闭 完满 域；不过对某些结果而言，在 $*$ 上「绝对地」工作至关重要。

对 $S\in\Perf_{\F_q}$，把
\[
X_S=Y_S/\varphi^\Z
\]
记作与 $S$ 相伴的\emph{相对曲线}。这里
\begin{equation}
Y_S=\Spa(W_{\OO_E}(R^+),W_{\OO_E}(R^+))\setminus V(\pi[\varpi]), \tag{2.1}
\end{equation}
其中 $S=\Spa(R,R^+)$ 是仿射 完满，$\varpi$ 是 $R$ 中的伪单值化元；按定义 $W_{\OO_E}(R^+)=R^+\hat\otimes_{\F_q}\OO_E\cong R^+[[\pi]]$（当 $E$ 为等特征时）。$Y_S$ 上 Frobenius $\varphi$ 的作用由分歧 Witt 向量 $W_{\OO_E}(R^+)$ 的 Frobenius 诱导。为说明对任意 $S$ 如何构造 $Y_S$ 及其 $\varphi^\Z$-作用，我们需要回忆关于菱形的一些基本事实。

\ysrc{127}回忆：若 $\mathcal F$ 是 完满 空间上的函子，把
\[
\mathcal F^{\diamondsuit}
\]
记作 $\F_p$-完满 空间上的函子，它把 $T\in\Perf_{\F_p}$ 送到
\[
\{(T^\sharp,\iota,s)\}/\sim,
\]
其中
\begin{itemize}[leftmargin=1.6em,itemsep=0.2em]
\item $T^\sharp$ 是 完满 空间；
\item $\iota:T\xrightarrow{\ \sim\ }(T^\sharp)^\flat$；
\item $s\in\mathcal F(T^\sharp)$；
\item 等价关系 $\sim$ 由下式给出：当且仅当存在同构 $f:T^\sharp\xrightarrow{\ \sim\ }T^{\sharp\prime}$ 满足
\begin{itemize}[leftmargin=1.2em,itemsep=0.1em]
\item $f^\flat\circ\iota=\iota'$，
\item $f^*s'=s$，
\end{itemize}
时，有 $(T^\sharp,\iota,s)\sim(T^{\sharp\prime},\iota',s')$。
\end{itemize}

对应 $\mathcal F\mapsto\mathcal F^{\diamondsuit}$ 把 $\Perf$ 上的 v-层送到 $\Perf_{\F_p}$ 上的 v-层；解析 adic 空间所表示的 v-层则被送到某一类特殊的 v-层，称为\emph{菱形}。菱形是 $\Perf_{\F_p}$ 上的 v-层，它们是「pro-étale 拓扑下的代数空间」，即可以写成
\[
X/R
\]
其中
\begin{itemize}[leftmargin=1.6em,itemsep=0.2em]
\item $X$ 是 $\F_p$-完满 空间；
\item $R\subset X\times_{\Spa(\F_p)}X$ 是由 完满 空间表示的等价关系，且两个态射 $R\rightrightarrows X$ 都是 pro-étale 的。
\end{itemize}
例如：
\begin{itemize}[leftmargin=1.6em,itemsep=0.2em]
\item 若 $X$ 是 完满 空间，则 $X^{\diamondsuit}$ 由 完满 空间 $X^\flat$ 表示，从而是菱形；
\item v-层 $\Spa(E)^{\diamondsuit}$ 是菱形，这从公式
\[
\Spa(E)^{\diamondsuit}=\Spa\big(\widehat{\overline E}^{\flat}\big)/\mathrm{Gal}(\overline E|E)
\]
可以看出。
\end{itemize}

\begin{rem}{2.1}
既然我们用了记号 $\mathrm{Gal}(\overline E|E)$，这里说明：对任一拓扑空间 $T$，函子
\[
\underline T:X\longmapsto C(|X|,T)
\]
是 完满 空间上的 v-层，即取值于 $T$ 的常预层所伴随的层。这一构造通常使我们能赋予「pro-étale $H$-挠子」这一概念确切含义，其中 $H$ 是局部投射有限的拓扑群：只需谈论 完满 空间上关于 pro-étale 拓扑的 $H$-挠子。例如，设 $\widehat{\mathbb G}_m$ 是 $\Z_p$ 上的乘法形式群，$\widehat{\mathbb G}_{m,\Q_p}$ 是它在 $\Spa(\Q_p)$ 上的刚性解析一般纤维。对数 $\log:\widehat{\mathbb G}_{m,\Q_p}\to\mathbb G_{a,\Q_p}$ 诱导一个 $\Z_p(1)$-pro-étale 挠子
\[
\varprojlim_{n\ge1}(\widehat{\mathbb G}_{m,\Q_p})^{\diamondsuit}\longrightarrow\mathbb G_{a,\Q_p}^{\diamondsuit},
\]
\label{rem:a3-2-1}
\end{rem}

\ysrc{128}其中
\begin{itemize}[leftmargin=1.6em,itemsep=0.2em]
\item 上述反向极限中的过渡态射由 $z\mapsto z^{\frac{m}{n}}$ 给出（当 $n|m$ 时）；
\item $\Z_p(1)$ 是 $\Z_p$ 在 $\Spa(\Q_p)^{\diamondsuit}$ 上的扭形式，它在 $\Spa(C_p^\flat)$ 上同构于 $\Z_p$，并由分圆特征给出。
\end{itemize}

\begin{rem}{2.2}
函子 $(-)^{\diamondsuit}$ 从 $\Perf$ 上的 v-层到 $\Perf_{\F_p}$ 上的 v-层，是倾斜的一种形式，它推广了 完满 空间的倾斜过程。更确切地说，完满 空间上的 v-拓扑是次典范的，并且有嵌入
\[
\Perf\subset\widetilde{\Perf}
\]
（v-层的拓扑斯）。于是下图交换
\begin{center}
\begin{tikzcd}[column sep=large,row sep=large]
\Perf \arrow[r, hook] \arrow[d, "(-)^\flat"'] & \widetilde{\Perf} \arrow[d, "(-)^{\diamondsuit}"] \\
\Perf_{\F_p} \arrow[r, hook] & \widetilde{\Perf}_{\F_p}
\end{tikzcd}
\end{center}
完满 空间的倾斜等价是指：若 $X$ 是 完满 空间，则
\[
(-)^\flat:\Perf_X\xrightarrow{\ \sim\ }\Perf_{X^\flat}.
\]
对 $\Perf$ 上的 v-层 $\mathcal F$，这延拓为拓扑斯的等价
\[
(-)^{\diamondsuit}:\widetilde{\Perf}_{/\mathcal F}\xrightarrow{\ \sim\ }\widetilde{\Perf}_{\F_p/\mathcal F^{\diamondsuit}}
\]
（局部化的拓扑斯）。典型地，这使我们能把 $\Q_p$-完满 空间上 v-层的拓扑斯看作一个特征 $p$ 的拓扑斯，
\[
(-)^{\diamondsuit}:\widetilde{\Perf}_{\Q_p}\xrightarrow{\ \sim\ }\widetilde{\Perf}_{\F_p/\Spa(\Q_p)^{\diamondsuit}}.
\]
\label{rem:a3-2-2}
\end{rem}

回到对任意 $S\in\Perf_{\F_q}$ 的 $X_S$ 的定义。对任意 $S\in\Perf_{\F_q}$，$Y_S$ 连同其 $\varphi$-作用是 $\Spa(E)$ 上唯一满足下述条件的解析 adic 空间：
\begin{enumerate}[label=(\arabic*),leftmargin=2.2em,itemsep=0.1em]
\item 其菱形由公式
\begin{equation}
Y_S^{\diamondsuit}=S\times_{\Spa(\F_q)^{\diamondsuit}}\Spa(E)^{\diamondsuit} \tag{2.2}
\end{equation}
给出，其中 $\varphi$ 作用为 $\Frob_S\times\mathrm{Id}$，$\Frob_S$ 是 $S$ 的绝对 $q$-Frobenius；
\item 对 $S$ 的任一仿射 完满 开子集，它由公式 (2.1) 给出。
\end{enumerate}
也就是说，公式 (2.2) 表明：对仿射 完满 的 $S$ 所作的构造，经由恒同 $|Y_S^{\diamondsuit}|=|Y_S|$ 粘成 adic 空间。事实上，正如「经典代数空间」的情形，可以这样定义菱形的拓扑空间：若 $X$ 是解析 adic 空间，则 $|X^{\diamondsuit}|=|X|$。

构造 $S\mapsto X_S$ 对 $S\in\Perf_{\F_q}$ 是函子性的。事实上我们有一族「经典曲线」，即 [8] 中所研究曲线的 adic 版本
\[
(X_{k(s),k(s)^+})_{s\in S},
\]
而 adic 空间 $X_S$ 是把这一族收集起来、由之构造出一族的方式。空间 $Y_S$ 与 $X_S$ 是 sous-完满 的 $E$-adic 空间，它们底变到 $\Spa(\breve E)$ 后成为 完满。

\ysrc{129}\subsubsection{叠 $\Bun_G$（[9, 第 III.1 章]）}

对 $S\in\Perf_{\F_q}$，把 $\Bun_G(S)$ 定义为 $X_S$ 上 $G$-丛的群胚。由于 $X_S$ 是 sous-完满 的，$X_S$ 上的向量丛有良好的概念；这里 $G$-丛按定义是从 $\Rep_E(G)$ 到 $\{X_S$ 上的向量丛$\}$ 的正合张量函子。

注意：当 $S=\Spa(R,R^+)$ 是仿射 完满 时，有一个相伴的「代数曲线」$\mathfrak X_S$，即 [8] 中那样的 $E$-概形（见 [9, 第 II.2.3 章]），并附带环化空间的态射 $\mathfrak X_S\to X_S$，它诱导 GAGA 等价（[9, 命题 II.2.7]）
\[
\{\mathfrak X_S\ \text{上的向量丛}\}\xrightarrow{\ \sim\ }\{X_S\ \text{上的向量丛}\}.
\]
于是 $X_S$ 上的 $G$-丛等同于概形 $\mathfrak X_S$ 上的 étale $G$-挠子。

第一个基本结果说：曲线上的 $G$-丛满足 v-拓扑的下降。这就是下述陈述。

\begin{thm}{2.3}
对应 $S\mapsto\Bun_G(S)$ 定义一个 v-叠
\[
\Bun_G\longrightarrow *.
\]
\label{thm:a3-2-3}
\end{thm}

利用 $X_S\hat\otimes_E\breve E$ 是 完满 的，这实际上容易归约到如下事实：完满 空间上的向量丛满足 v-拓扑的下降（[25, 引理 17.1.8]）。

\subsection{点（[9, 第 II.2.1 章]）}

回忆下述构造。设 $S\in\Perf_{\F_q}$。有函子
\[
\mathrm{Isocrystals}\longrightarrow\{X_S\ \text{上的向量丛}\},\qquad (D,\varphi)\longmapsto\mathcal E(D,\varphi),
\]
由公式
\[
\mathcal E(D,\varphi)=Y_S\overset{\varphi^\Z}{\times}D\longrightarrow Y_S/\varphi^\Z=X_S
\]
给出。这里等晶体是有限维 $\breve E$-向量空间 $D$ 连同 $\sigma$-半线性自同构 $\varphi:D\xrightarrow{\sim}D$。我们用到了态射 $Y_S\to\Spa(\breve E)$ 与 $Y_S$ 上的 $\varphi$-作用、$\breve E$ 上的 $\sigma$-作用相容。上述构造 $(D,\varphi)\mapsto\mathcal E(D,\varphi)$ 升级为映射
\[
G(\breve E)\longrightarrow\{X_S\ \text{上的 }G\text{-丛}\},\qquad b\longmapsto\mathcal E_b.
\]
回忆 Kottwitz 集是
\[
B(G)=G(\breve E)/\sigma\text{-共轭},\qquad b\sim gbg^{-\sigma},
\]
见 [16]。现在有下述结果。

\begin{thm}{2.4}（[6],[1]）
当 $S=\Spa(F,F^+)$ 是几何点，即 $F$ 是代数闭 完满 域时，
\[
B(G)\xrightarrow{\ \sim\ }\{X_{F,F^+}\ \text{上的 }G\text{-丛}\}/\sim,\qquad [b]\longmapsto[\mathcal E_b].
\]
\label{thm:a3-2-4}
\end{thm}

\ysrc{130}在 $\mathrm{GL}_n$ 情形，这个结果是 [8] 的主要结果之一：代数曲线 $\mathfrak X_F$ 上的任一向量丛同构于 $\bigoplus_i\OO(\lambda_i)$，其中斜率 $\lambda_i\in\Q$。我们在 [9, 第 II.2.4 章] 中用菱形理论给出这个结果的新证明。[8] 中的证明依赖于 Lubin–Tate 空间与 Drinfeld 空间的周期域。

\begin{rem}{2.5}
在上述对应 $(D,\varphi)\mapsto\mathcal E(D,\varphi)$ 中，斜率是反过来的：若 $(D,\varphi)$ 是斜率为 $\lambda$ 的等斜等晶体，则 $\mathcal E(D,\varphi)$ 是半稳定的，其 Harder–Narasimhan 斜率为 $-\lambda$。若 $(D,\varphi)$ 是斜率为 $\lambda$ 的单等晶体，则 $\mathcal E(D,\varphi)\simeq\OO(-\lambda)$。
\label{rem:a3-2-5}
\end{rem}

于是我们得到集合的恒同
\[
B(G)=|\Bun_G|.
\]
这个等式是研究 $\Bun_G$ 几何的起点。

\subsection{Banach–Colmez 空间上的结果（[9, 第 II.2 章], [18]）}

Banach–Colmez 空间在某种意义上就是菱形理论中的线性对象。它们是我们到处用来做几何的基本构件。它们来自曲线上向量丛的相对层上同调。下文我们谈论\emph{局部空间菱形}。事实上菱形的范畴太大；为得到与解析 adic 空间范畴性质相似的范畴，需要加以限制。按定义，菱形 $X$ 是\emph{空间的}，若
\begin{enumerate}[label=(\arabic*),leftmargin=2.2em,itemsep=0.1em]
\item 它是拟紧拟分离的；
\item 其相伴拓扑空间是谱空间。
\end{enumerate}
典型地，若 $X$ 是解析 adic 空间，则 $X^{\diamondsuit}$ 是局部空间的。局部空间菱形的性质有两个例子：
\begin{enumerate}[label=(\arabic*),leftmargin=2.2em,itemsep=0.1em]
\item 若 $f:X\to Y$ 是局部空间菱形的态射，则 $|f|:|X|\to|Y|$ 是泛化的；
\item 若 $X$ 是空间菱形，则 $|X|$ 的任一 pro-可构造泛化子集 $Z$ 定义空间菱形 $i:Y\hookrightarrow X$，使 $i$ 拟紧且 $|Y|=Z$。
\end{enumerate}
第一条性质推广了解析 adic 空间的通常性质。这里是第二条的一个应用：若 $X$ 是拟紧拟分离的解析 adic 空间，则 $|X|=|X^{\diamondsuit}|$ 的任一 pro-可构造泛化子集定义 $X^{\diamondsuit}$ 的空间子菱形，而后者一般可能不与任何解析 adic 空间相伴。关于这一点更多内容见 [23]。

\subsubsection{在一个点上}

设 $C|\F_q$ 是代数闭 完满 域。关于曲线上向量丛的上同调层，我们有下述结果。这里所谓斜率指 Harder–Narasimhan 斜率；所谓层指 $\Perf_C$ 上的 v-层。当我们说态射是 $\ell$-上同调光滑的，意思是在 Verdier 意义下「相对 Poincaré 对偶成立」。

\begin{thm}{2.6}
设 $\mathcal E$ 是 $X_C$ 上的向量丛。
\begin{enumerate}[label=(\arabic*),leftmargin=2.2em,itemsep=0.2em]
\item 若 $\mathcal E$ 的斜率 $>0$，则当 $S$ 是仿射 完满 时 $H^1(X_S,\mathcal E|_{X_S})=0$，且 $S\mapsto H^0(X_S,\mathcal E|_{X_S})$ 是层。这个层
\[
\mathcal{BC}(\mathcal E)\longrightarrow\Spa(C)
\]
是分离的 $\ell$-上同调光滑局部空间菱形，维数为 $\deg(\mathcal E)$。
\ysrc{131}\item 若进一步 $\mathcal E$ 的斜率在 $]0,[E:\Q_p]]$ 中，则它由一个开的 完满 球表示，更确切地说，是 $\F_q$ 上一个形式 $p$-可除群的全覆盖。
\item 若 $\mathcal E$ 的斜率 $<0$，则对任意 $S$ 有 $H^0(X_S,\mathcal E|_{X_S})=0$，且 $S\mapsto H^1(X_S,\mathcal E)$ 是层。这个层
\[
\mathcal{BC}(\mathcal E[1])\longrightarrow\Spa(C)
\]
是分离的 $\ell$-上同调光滑局部空间菱形，维数为 $-\deg(\mathcal E)$。
\end{enumerate}
\label{thm:a3-2-6}
\end{thm}

我们解释一些特殊情形。设 $E=\Q_p$。若 $\mathcal G$ 是 $\F_q$ 上的形式 $p$-可除群，其共变 Dieudonné 模为 $(D,\varphi)$，则
\[
\mathcal{BC}(\mathcal E(D,\varphi)(1))\simeq\mathcal G\times_{\Spa(\F_q)}\Spa(C),
\]
其中在此公式中 $\mathcal G\simeq\Spf(\F_q[[x_1,\cdots,x_d]])$ 是形式概形，我们把它看作 $\Perf_{\F_q}$ 上的 v-层。为此回忆：若 $\mathfrak X$ 是 $\F_p$-形式概形，则相伴的 v-层是函子 $(R,R^+)\mapsto\mathfrak X(R^+)$ 在 $\F_p$-仿射 完满 代数上关于解析拓扑的层化（这一解析层化是 v-层；理由是：把它底变到任一仿射 完满 $\Spa(R,R^+)$ 后，它可由 完满 空间表示）。在 $\Spf(\F_q[[x_1,\cdots,x_d]])$ 的情形，无需关于解析拓扑作这样的层化，它就是函子 $(R,R^+)\mapsto(R^{\circ\circ})^d$。

对一般的 $E$ 与 $>0$ 的斜率，可以找到正合列
\[
0\longrightarrow\mathcal E''\longrightarrow\mathcal E'\longrightarrow\mathcal E\longrightarrow 0,
\]
其中 $\mathcal E''$ 是平凡向量丛，$\mathcal E'$ 的斜率在 $]0,1]$ 中。于是 $\mathcal{BC}(\mathcal E)$ 同构于 $\mathcal{BC}(\mathcal E')$ 对一个 pro-étale 等价关系的商；该等价关系由 $V$ 在这个开球上的自由作用给出，其中 $V=H^0(X_C,\mathcal E'')$ 是有限维 $E$-向量空间。

在负斜率情形，固定 $C$ 到 $E$ 的一个直到化 $C^\sharp$ 之后，$X_C$ 上有一个闭点 $\infty$，其剩余域为 $C^\sharp$。固定恒同 $\OO_{X_C}(1)\simeq\OO_{X_C}(\infty)$ 后，我们得到同构
\[
\mathcal{BC}(\OO(-1)[1])\simeq\mathbb G^{\diamondsuit}_{a,C^\sharp/E},
\]
它显然是分离的 $\ell$-上同调光滑菱形。

\subsubsection{绝对版本（[9, 第 II.2.2 章]）}

下述结果对负 HN 斜率是新的，对正斜率则不然（见 [7]）。曲线上向量丛的相对上同调在 $*$ 上绝对地不能由菱形表示，但打洞之后可以。

\begin{thm}{2.7}
设 $(D,\varphi)$ 是等晶体。
\begin{enumerate}[label=(\arabic*),leftmargin=2.2em,itemsep=0.2em]
\item 若 $(D,\varphi)$ 的斜率 $<0$，则把 $S$ 送到 $H^0(X_S,\mathcal E(D,\varphi))$ 的 v-层
\[
\mathcal{BC}(\mathcal E(D,\varphi))\longrightarrow *
\]
满足：$\mathcal{BC}(\mathcal E(D,\varphi))\setminus\{0\}$ 是 $*$ 上的空间菱形。
\item 若 $(D,\varphi)$ 的斜率 $>0$，则把 $S$ 送到 $H^1(X_S,\mathcal E(D,\varphi))$ 的 v-层
\[
\mathcal{BC}(\mathcal E(D,\varphi)[1])\longrightarrow *
\]
满足：$\mathcal{BC}(\mathcal E(D,\varphi)[1])\setminus\{0\}$ 是 $*$ 上的空间菱形。
\end{enumerate}
\label{thm:a3-2-7}
\end{thm}

\ysrc{132}最困难的是第 (2) 点：这里必须用空间菱形的 Artin 判别法（[23, 定理 12.18]），先证明 $\mathcal{BC}(\mathcal E(D,\varphi)[1])$ 是空间 v-层，再给出它被局部闭空间菱形的分层。

\begin{rem}{2.8}
我们之所以把 [23, 定理 12.18] 称为空间菱形的 Artin 判别法，理由如下。通常关于 fppf 层 $X$ 是代数空间的 Artin 判别法，涉及一些整体假设（如某种有限表现假设），再配合某个局部假设（涉及 $X$ 在 $|X|$ 中某点处的形式完备化）。[23] 的定理 2.18 作用方式相同：$\F_p$-完满 空间上的小 v-层 $X$ 是空间菱形，若
\begin{enumerate}[label=(\arabic*),leftmargin=2.2em,itemsep=0.1em]
\item 它是空间 v-层，即拟紧拟分离，且 $|X|$ 是谱空间；
\item 对任意 $x\in|X|$，局部化 $X_x=\varprojlim_{U\ni x}U$（$U$ 取遍 $x$ 的邻域）是菱形。
\end{enumerate}
\label{rem:a3-2-8}
\end{rem}

在这个判别法的所有应用中（定理 2.7，或 $B_{dR}$-仿射 Grassmann 簇中闭 Schubert 胞腔是空间菱形这一事实的证明），第 (2) 点都通过给出 $|X|$ 被局部闭泛化子集的一个分层、而每个这样的子集可由菱形表示来得到。我们在定理 2.7 中证明 $X$ 是空间 v-层的方式，是构造一个 v-满的 $\ell$-上同调光滑态射 $\widetilde X\to X$，其中 $\widetilde X$ 是空间菱形，并利用
\begin{enumerate}[label=(\arabic*),leftmargin=2.2em,itemsep=0.1em]
\item $\ell$-上同调光滑态射是开的；
\item 若 $Z$ 是谱空间，$R\subset Z\times Z$ 是 pro-可构造等价关系，且两个映射 $R\rightrightarrows Z$ 都是开的，则 $Z/R$ 是谱空间。
\end{enumerate}

\begin{exam}{2.9}
最简单的例子是 $(D,\varphi)$ 的斜率在 $[-[E:\Q_p],0[$ 中的情形。此时
\[
\mathcal{BC}(\mathcal E(D,\varphi))\simeq\Spa\big(\overline{\F}_q[[x_1^{1/p^\infty},\cdots,x_d^{1/p^\infty}]],\overline{\F}_q[[x_1^{1/p^\infty},\cdots,x_d^{1/p^\infty}]]\big)
\]
它不是 完满 空间（不是解析 adic 空间），也不是菱形；但去掉点 $V(x_1,\dots,x_d)$ 之后，它成为拟紧 完满 空间。\tn{上式照录底本：底本把 $\Spa(R,R^+)$ 的两个参数都印作 $\overline{\F}_q[[x_1^{1/p^\infty},\cdots,x_d^{1/p^\infty}]]$。按 $\Spa$ 的惯例，第一个参数应取分式域 $\overline{\F}_q((x_1^{1/p^\infty},\cdots,x_d^{1/p^\infty}))$，第二个参数取它的整环。}
\label{exam:a3-2-9}
\end{exam}

我们更一般地解释正 HN 斜率情形，$\mathcal{BC}(\OO(d))\setminus\{0\}$，$d\in\N_{\ge1}$。对 $d=1$，它同构于 $\Spa(\F_q((x^{1/p^\infty})))$。一般地令 $\Delta=\{(\lambda_1,\dots,\lambda_d)\in(E^\times)^d\mid\prod_i\lambda_i=1\}$。则可证明乘积态射诱导 pro-étale 商同构
\[
\big(\mathcal{BC}(\OO(1))\setminus\{0\}\times\cdots\times\mathcal{BC}(\OO(1))\setminus\{0\}\big)/\Delta\rtimes\mathfrak S_d\xrightarrow{\ \sim\ }\mathcal{BC}(\OO(d))\setminus\{0\},
\]
其中 $\Delta\rtimes\mathfrak S_d$ 的作用不自由。无论如何，这足以证明 $\mathcal{BC}(\OO(d))\setminus\{0\}$ 是空间菱形。细节见 [7]。

我们解释负 HN 斜率情形 $\mathcal{BC}(\OO(-1)[1])\setminus\{0\}$。它分类 $X_S$ 上 $\OO(1)$ 被 $\OO$ 的扩张 $0\to\OO\to\mathcal E\to\OO(1)\to0$，这些扩张在 $S$ 上逐纤维非分裂。于是，几何地逐纤维于 $S$，$\mathcal E\simeq\OO(\frac12)$（这是定理 2.4 的特例）。

\ysrc{133}由此推出
\[
\mathcal{BC}(\OO(-1)[1])\setminus\{0\}\simeq\mathcal{BC}(\OO(\tfrac12))\setminus\{0\}/(D^\times)^1,
\]
其中 $D$ 是 $E$ 上的四元数代数。这使人可以验证它是空间菱形。

这些绝对 Banach–Colmez 空间在我们的工作中起关键作用，见 §2.8 与空间菱形 $\widetilde{\mathcal M}_b$。

\subsubsection{Banach–Colmez 空间族}

在 [9] 中我们用菱形理论重做了 [14] 的结果。由此得到下述关于曲线上向量丛族的相对上同调、从而 Banach–Colmez 空间族的结果。

\begin{thm}{2.10}
设 $S\in\Perf_{\F_q}$，$\mathcal E$ 是 $X_S$ 上的向量丛。
\begin{enumerate}[label=(\arabic*),leftmargin=2.2em,itemsep=0.2em]
\item v-层
\[
T/S\mapsto H^0(X_T,\mathcal E|_{X_T})
\]
是局部空间菱形。
\item 若在 $S$ 上逐纤维 $\mathcal E$ 的斜率 $>0$，则与 $T/S\mapsto H^1(X_T,\mathcal E|_{X_T})$ 相伴的层为零。此外 $T/S\mapsto H^0(X_T,\mathcal E|_{X_T})$ 是 $S$ 上分离的 $\ell$-上同调光滑局部空间菱形。
\item 若在 $S$ 上逐纤维 $\mathcal E$ 的斜率 $<0$，则对任意 $T/S$ 有 $H^0(X_T,\mathcal E|_{X_T})=0$，且 $T/S\mapsto H^1(X_T,\mathcal E|_{X_T})$ 是层。这个层还是 $S$ 上分离的 $\ell$-上同调光滑局部空间菱形。
\end{enumerate}
\label{thm:a3-2-10}
\end{thm}

这给我们提供了在局部空间菱形范畴中构造一些良好线性对象的方法。这些就是我们使用的基本几何对象。

\subsection{几何结构（[9, 第 IV.1 章]）}

v-叠 $\Bun_G$ 先验地是抽象对象，但下述结果说它有良好的几何结构。

\begin{thm}{2.11}
下列结论成立。
\begin{enumerate}[label=(\arabic*),leftmargin=2.2em,itemsep=0.2em]
\item v-叠 $\Bun_G$ 是 Artin v-叠：
\begin{enumerate}[label=(\alph*),leftmargin=1.8em,itemsep=0.1em]
\item $\Bun_G$ 的对角线可由局部空间菱形表示；
\item 存在局部空间菱形 $U$ 以及分离的、满的 $\ell$-上同调光滑态射 $U\to\Bun_G$。
\end{enumerate}
\item $\Bun_G\to*$ 是维数为 $0$ 的分离 $\ell$-上同调光滑态射。其对偶复形同构于 $\Lambda$。
\end{enumerate}
\label{thm:a3-2-11}
\end{thm}

\begin{rem}{2.12}
必须注意：与光滑射影曲线上向量丛的「经典叠」相反，$\Bun_G$ 的对角态射不是拟紧的。特别地，若 $U\subset\Bun_G$ 是拟紧开子集，则包含 $j:U\hookrightarrow\Bun_G$ 一般不是拟紧的，典型地不能应用拟紧底变（[23, 推论 16.10]）。一般地 $R^ij_*\overline{\F}_\ell$ 对 $i>0$ 非零，这与拟紧拟分离开浸入的情形相反。
\label{rem:a3-2-12}
\end{rem}

前述定理中的第 (1)(a) 点容易由定理 2.10 推出。典型地，对 $\mathrm{GL}_n$，若 $\mathcal E_1,\mathcal E_2$ 是 $X_S$ 上的向量丛，则 v-层 $\mathrm{Isom}(\mathcal E_1,\mathcal E_2)\to S$ 是 $\mathcal{BC}(\mathcal E_1^\vee\otimes\mathcal E_2)$ 中的开子集。

\ysrc{134}证明该定理中其他结果的一种办法，是使用所谓的 Beauville–Laszlo 单值化（[9, 第 III.3 章]）。更确切地说，设
\[
\mathrm{Gr}_G\longrightarrow\mathrm{Spd}(\breve E)
\]
是 $B_{dR}$-仿射 Grassmann 簇（[25, 第 19 讲]）。若 $S$ 是 $\F_q$-完满 空间并带有态射 $S\to\mathrm{Spd}(\breve E)$，即 $\breve E$ 上的一个直到化 $S^\sharp$，则
\[
S^\sharp\hookrightarrow X_S
\]
是 Cartier 除子。现在 $\mathrm{Gr}_G(S)$ 由这样的直到化 $S^\sharp$、$X_S$ 上的 $G$-丛 $\mathcal E$，以及同构
\[
\mathcal E_1|_{X_S\setminus S^\sharp}\xrightarrow{\ \sim\ }\mathcal E|_{X_S\setminus S^\sharp}
\]
组成，其中该同构「沿除子 $S^\sharp$ 亚纯」，$\mathcal E_1$ 是平凡 $G$-丛。这里亚纯的定义可归约到向量丛情形：要求对 $G$ 的任一线性表示 $\rho$，向量丛的同构 $\rho_*\mathcal E_1|_{X_S\setminus S^\sharp}\xrightarrow{\sim}\rho_*\mathcal E|_{X_S\setminus S^\sharp}$ 是亚纯的。在向量丛情形，即 $G=\mathrm{GL}_n$（$n\ge1$）时，该同构亚纯是指：在 $S$ 上局部地，它来自某个态射
\[
\mathcal E_1\longrightarrow\mathcal E(kD)
\]
到 $X_S\setminus S^\sharp$ 的限制，其中 $k\gg0$，$D$ 是由 $S^\sharp\hookrightarrow X_S$ 定义的 Cartier 除子。

此外回忆 $\mathrm{Gr}_G$ 可由 ind-菱形表示。为把它写成 ind-菱形，设 $G$ 分裂以简化记号，并令 $T\subset B$ 是 Borel 子群中的极大环面。对每个 $\mu\in X_*(T)^+$，有相应的 Schubert 胞腔，包含在闭 Schubert 胞腔
\[
\mathrm{Gr}_{G,\mu}\subset\mathrm{Gr}_{G,\le\mu}
\]
之中；其中开 Schubert 胞腔是 $\mathrm{Spd}(\breve E)$ 上维数为 $\langle\mu,2\rho\rangle$ 的分离 $\ell$-上同调光滑局部空间菱形，闭胞腔是 $\mathrm{Spd}(\breve E)$ 上的真空间菱形。于是
\[
\mathrm{Gr}_G=\varinjlim_{\mu\in X_*(T)^+}\mathrm{Gr}_{G,\le\mu},
\]
其中按定义 $\mu_1\le\mu_2$ 指 $\mu_2-\mu_1\in\N.\check\Phi$，即余根的正和。然后我们证明下述结果。

\begin{thm}{2.13}（[9, 命题 III.3.1, 定理 IV.1.19]）
下列结论成立：
\begin{enumerate}[label=(\arabic*),leftmargin=2.2em,itemsep=0.2em]
\item Beauville–Laszlo 态射 $\mathrm{Gr}_G\to\Bun_G$ 是 v-满的。
\item 态射
\[
\coprod_{\mu\in X_*(T)^+}\big[G(E)\backslash\mathrm{Gr}_{G,\mu}\big]\longrightarrow\Bun_G
\]
是分离 $\ell$-上同调光滑的，并给出 $\Bun_G$ 的一个表现，即把它表示成 $*$ 上分离 $\ell$-上同调光滑局部空间菱形。
\end{enumerate}
\label{thm:a3-2-13}
\end{thm}

这张 $\Bun_G$ 的图卡容易构造。然而它不是我们用来分析 $\Bun_G$ 上层的那一张，见 §2.8。

\ysrc{135}\subsection{HN 分层与连通分支（[9, 第 III.2.2、III.2.4、IV.1.2.2 章]）}

设 $G^*$ 是 $G$ 的拟分裂内型，$A\subset T\subset B$ 是 Borel 子群中极大环面内的极大分裂环面。有恒同
\[
X_*(A)^+_\Q=[\mathrm{Hom}(\mathbb D,G_{\overline E})/G(\overline E)]^\Gamma,
\]
其中 $\mathbb D$ 是斜率 pro-环面，$X^*(\mathbb D)=\Q$。我们称 $\pi_1(G)$ 为 Borovoi 基本群。

\begin{thm}{2.14}
下列结论成立。
\begin{enumerate}[label=(\arabic*),leftmargin=2.2em,itemsep=0.2em]
\item 由 Harder–Narasimhan 多边形给出的映射 $|\Bun_G|\to X_*(A)^+_\Q$ 是上半连续的。
\item 由第一 Chern 类给出的映射 $|\Bun_G|\to\pi_1(G)_\Gamma$ 局部常值，且纤维连通。
\end{enumerate}
\label{thm:a3-2-14}
\end{thm}

用恒同 $B(G)=|\Bun_G|$ 来说，$\mathcal E_b$ 的 Harder–Narasimhan 多边形是 $w_0.(-\nu_b)$，其中 $\nu_b\in X_*(A)^+_\Q$ 是 Newton 点，$w_0$ 是 Weyl 群中长度最大的元素。$\mathcal E_b$ 的第一 Chern 类是 $-\kappa(b)$，其中 $\kappa$ 是 Kottwitz 映射 $\kappa:B(G)\to\pi_1(G)_\Gamma$。

\begin{rem}{2.15}
当导子群 $G_{der}$ 单连通时，$|\Bun_G|\to\pi_1(G)_\Gamma$ 的局部常值性容易看出：它实际上就由 $|\Bun_G|\to|\Bun_G/G_{der}|$ 给出。困难的情形是 $G_{der}$ 不单连通，[9, 第 III.2.4 章] 处理了这一情形。
\label{rem:a3-2-15}
\end{rem}

为证明 Harder–Narasimhan 多边形的上半连续性，我们给出一个用菱形理论的证明。我们解释 $\mathrm{GL}_n$ 情形。设 $\mathcal E$ 是 $X_S$ 上的向量丛。$\mathcal E$ 在 $S$ 的几何点处的 HN 多边形首斜率 $\ge\lambda$，当且仅当在该几何点处存在非零态射 $\OO(\lambda)\to\mathcal E$。从 $\OO(\lambda)$ 到 $\mathcal E$ 的非零态射的模叠是 $\mathcal{BC}(\OO(-\lambda)\otimes\mathcal E)\setminus\{0\}\to S$。现在利用态射
\[
\mathcal{BC}(\OO(-\lambda)\otimes\mathcal E)\setminus\{0\}/\pi^\Z\longrightarrow S
\]
是局部空间菱形的真态射。该态射在 $S$ 中的像，即 $\mathcal E$ 的多边形首斜率 $\ge\lambda$ 的轨迹，从而是闭的。把这一论证用于 $\mathcal E$ 的外幂，即可得出 $\mathcal E$ 的 HN 多边形的上半连续性。

最后我们指出，特别地，半稳定轨迹
\[
\Bun_G^{ss}\subset\Bun_G
\]
是开的。我们将在后面详细描述这个开子叠的结构。

\begin{thm}{2.16}（[26]）
$|\Bun_G|$ 的拓扑是由嵌入 $B(G)\hookrightarrow X_*(A)^+_\Q\times\pi_1(G)_\Gamma$ 以及 $X_*(A)^+_\Q$ 上的序诱导的。
\label{thm:a3-2-16}
\end{thm}

对 $\mathrm{GL}_n$，集合 $B(G)$ 可描述为一组 Newton 多边形。此时结论是：$|\Bun_G|$ 的点 $x$ 是 $y$ 的特化，当且仅当与 $x$ 相伴的多边形位于与 $y$ 相伴的多边形之上，且两者端点相同。这一特殊情形在 [4] 中处理。

\ysrc{136}\subsection{HN 分层作为分类叠}

\subsubsection{半稳定轨迹（[9, 第 III.4 章]）}

Kottwitz $\kappa$ 映射诱导双射
\[
B(G)_{\mathrm{basic}}\xrightarrow{\ \sim\ }\pi_1(G)_\Gamma.
\]
于是每个连通分支中恰有一个半稳定点。现在有下述结果。对基本的 $b$，层
\[
S\mapsto\mathrm{Aut}(\mathcal E_b/X_S)
\]
就是层 $\underline{G_b(E)}$，它在 $S$ 上的取值是连续函数 $|S|\to G_b(E)$。这里 $G_b$ 是 $G$ 的内型，即所谓的 $G$ 的扩充纯内型。例如 $G_1=G$。于是有下述结果。

\begin{thm}{2.17}
与基本的 $b$ 相伴的分层等同于 pro-étale $G_b(E)$-挠子的分类叠
\[
[*/G_b(E)].
\]
\label{thm:a3-2-17}
\end{thm}

这里恒同把 $X_S$ 上、在 $S$ 上几何逐纤维同构于 $\mathcal E_b$ 的 $G$-丛 $\mathcal E$ 送到挠子
\[
T/S\longmapsto\mathrm{Isom}(\mathcal E_b,\mathcal E|_{X_T}).
\]

我们指出：利用 étale 分离态射满足 v-下降（[23, 命题 9.7]），$G_b(E)$-pro-étale 挠子等同于 $G_b(E)$-v-挠子。

\begin{exam}{2.18}
对线性群，上述结论说：$X_S$ 上几何逐纤维斜率 $0$ 半稳定的向量丛 $\mathcal E$，与几何茎为有限维的 $E$-向量空间的 pro-étale 局部常值层之间存在等价。这里对应把这样的 $\mathcal E$ 送到层 $T/S\mapsto H^0(X_T,\mathcal E|_{X_T})$；反方向把 $\mathcal F$ 送到 $\mathcal F\otimes_E\OO_{X_S}$。
\label{exam:a3-2-18}
\end{exam}

作为推论，我们得到分解
\[
\Bun_G^{ss}=\coprod_{[b]\ \mathrm{basic}}[*/G_b(E)].
\]

\subsubsection{更一般的分层（[9, 第 III.5 章]）}

固定 $[b]\in B(G)$。$\mathcal E_b$ 的自同构层的结构一般比基本情形更复杂，意思是其单位连通分支非平凡。更确切地说，有下述结果。

\begin{thm}{2.19}
把 $S$ 送到 $\mathrm{Aut}(\mathcal E_b/X_S)$ 的、$\mathcal E_b$ 的自同构层 $\widetilde G_b$ 具有形状
\[
\widetilde G_b=\widetilde G_b^0\rtimes\underline{G_b(E)},
\]
其中 $\widetilde G_b^0$ 是幂幺群菱形，是正 Banach–Colmez 空间的逐次扩张。还有 $\dim(\widetilde G_b^0)=\langle\nu_b,2\rho\rangle$。
\label{thm:a3-2-19}
\end{thm}
